\section*{Tag 13 — Datenplatzierung im Datamatrix}

\subsection*{Bleistift-Übung 1: Utah-Form für Codeword 1 (HONIG)}

Codeword 1 von HONIG ist \texttt{73}. In Binär: $73 = 64 + 8 + 1 =
\texttt{01001001}$.

\begin{center}
\small
\begin{tabular}{c | c c c}
\toprule
& Spalte $-2$ & Spalte $-1$ & Spalte $0$ \\
\midrule
Zeile $-2$ & Bit 1 = \textbf{0} & Bit 2 = \textbf{1} & --- \\
Zeile $-1$ & Bit 3 = \textbf{0} & Bit 4 = \textbf{0} & Bit 5 = \textbf{1} \\
Zeile $0$  & Bit 6 = \textbf{0} & Bit 7 = \textbf{0} & Bit 8 = \textbf{1} \\
\bottomrule
\end{tabular}
\end{center}

Der Cursor steht beim Algorithmus zu Beginn bei \texttt{(row, col) =
(4, 0)}. Die acht Modul-Positionen sind also:
$(2, -2), (2, -1), (3, -2), (3, -1), (3, 0), (4, -2), (4, -1), (4, 0)$.

Die negativen Spalten wickeln sich am Symbol-Rand auf die rechte
Seite — Position $(2, -2)$ wird beispielsweise zu $(2 + 0, -2 + 10) =
(2, 8)$ in der 10$\times$10-Innenfläche, also in Symbol-Koordinaten
(also +1 für L-Pattern-Offset) zu Modul $(3, 9)$.

(Wer das genauer nachrechnen will: die Wrap-Formel ist
\texttt{col $\to$ col + n}, \texttt{row $\to$ row + 4 - ((n+4) mod 8)},
was für $n=10$ ergibt: $4 - 6 = -2$ als Zeilen-Verschiebung. Die
genauen Koordinaten kannst du mit \texttt{\_place\_one} in Thonny
einzeln durchspielen.)

\subsection*{Bleistift-Übung 2: HONIG zeichnen}

Vorgehen:

\begin{enumerate}
\item Trage in der Hilfstabelle der Vorlage die zwölf Codewort-Bytes
  ein:
  \texttt{[73, 80, 79, 74, 72, 70, 186, 97, 167, 44, 40, 243]}
\item Pro Codeword die acht Bits (MSB links) ausrechnen.
\item Im Vorbild-Feld sind die Module bereits eingezeichnet —
  vergleiche jeden Modul mit deiner Tabelle.
\end{enumerate}

Wenn du Lust hast, das gleiche im leeren Feld nochmal selbst zu
zeichnen, ist das die beste Methode, den Algorithmus wirklich zu
internalisieren. Der erste Modul ist bereits durch das L-Pattern
gesetzt — du beginnst mit dem Datenbereich.

\subsection*{Bleistift-Übung 3: Eigenes Symbol zeichnen}

Aus Etappe 12 hast du die ECC für GRETA berechnet:
\texttt{[72, 83, 70, 85, 66, 64, 90, 71, 128, 186, 106, 125]}.

Hier die komplette Bit-Tabelle dazu:

\begin{center}
\small
\begin{tabular}{c c | c c c c c c c c}
\toprule
\textbf{Nr.} & \textbf{Wert} & \textbf{Bit 1} & \textbf{2} & \textbf{3} & \textbf{4} & \textbf{5} & \textbf{6} & \textbf{7} & \textbf{Bit 8} \\
                &           & (128) & (64) & (32) & (16) & (8) & (4) & (2) & (1) \\
\midrule
\phantom{0}1  &  72 & 0 & 1 & 0 & 0 & 1 & 0 & 0 & 0 \\
\phantom{0}2  &  83 & 0 & 1 & 0 & 1 & 0 & 0 & 1 & 1 \\
\phantom{0}3  &  70 & 0 & 1 & 0 & 0 & 0 & 1 & 1 & 0 \\
\phantom{0}4  &  85 & 0 & 1 & 0 & 1 & 0 & 1 & 0 & 1 \\
\phantom{0}5  &  66 & 0 & 1 & 0 & 0 & 0 & 0 & 1 & 0 \\
\phantom{0}6  &  64 & 0 & 1 & 0 & 0 & 0 & 0 & 0 & 0 \\
\phantom{0}7  &  90 & 0 & 1 & 0 & 1 & 1 & 0 & 1 & 0 \\
\phantom{0}8  &  71 & 0 & 1 & 0 & 0 & 0 & 1 & 1 & 1 \\
\phantom{0}9  & 128 & 1 & 0 & 0 & 0 & 0 & 0 & 0 & 0 \\
10            & 186 & 1 & 0 & 1 & 1 & 1 & 0 & 1 & 0 \\
11            & 106 & 0 & 1 & 1 & 0 & 1 & 0 & 1 & 0 \\
12            & 125 & 0 & 1 & 1 & 1 & 1 & 1 & 0 & 1 \\
\bottomrule
\end{tabular}
\end{center}

Trage diese zwölf Bytes ins zweite Werkstatt-Feld ein und zeichne
das Datenraster Modul für Modul. Wenn du fertig bist, vergleiche
mit der Python-Ausgabe — \texttt{place\_codewords([72, 83, 70, 85,
66, 64, 90, 71, 128, 186, 106, 125], 12)} sollte exakt deinem
gezeichneten Symbol entsprechen.

Tipp zur Selbstkontrolle: nimm das gerenderte PNG aus Etappe 14
(wenn der finale Encoder steht) und überlege es transparent mit
deiner Hand-Zeichnung — jeder Unterschied ist ein Hand-Fehler.
